% Generated draft: requires mathematical and source review.
\documentclass[11pt]{article}
\usepackage[margin=25mm]{geometry}
\usepackage{fontspec}
\setmainfont{Latin Modern Roman}
\usepackage{amsmath,amssymb,mathtools,unicode-math}
\setmathfont{Latin Modern Math}
\usepackage{enumitem,needspace,xurl,hyperref}
\setlength{\parindent}{0pt}
\setlength{\parskip}{5pt}
\setlength{\emergencystretch}{3em}
\newcommand{\N}{\mathbb{N}}
\newcommand{\Z}{\mathbb{Z}}
\newcommand{\Q}{\mathbb{Q}}
\newcommand{\R}{\mathbb{R}}
\newcommand{\C}{\mathbb{C}}
\newcommand{\F}{\mathbb{F}}
\newcommand{\A}{\mathbb{A}}
\newcommand{\PP}{\mathbb{P}}
\newcommand{\E}{\mathbb{E}}
\newcommand{\eps}{\varepsilon}
\newcommand{\dd}{\,\mathrm{d}}
\begin{document}
\section*{UnsolvedMath ID: 30002359}
\textbf{Source problem code:} catalogue entry
\subsection{Definitions and mathematical statement}
Let \(F\) be a field. A quadratic form over \(F\) is a map \(q\colon V\to F\) on a finite-dimensional \(F\)-vector space such that \(q(av)=a^2q(v)\) and \(b_q(v,w)=q(v+w)-q(v)-q(w)\) is bilinear. Its dimension is \(\dim_F V\). The form is anisotropic if \(q(v)=0\) implies \(v=0\), and universal if every element of \(F^\times\) is represented by \(q\). Under the polar-form convention, nonsingularity means that \(\{v\in V:b_q(v,w)=0\text{ for every }w\in V\}=\{0\}\); the intended source convention requires confirmation.

The supplied record uses \(u(F)\) and \(\widetilde u(F)\) without defining them. They are retained as source notation: their defining classes of forms, any additional restrictions, and supremum conventions must be recovered from the source before these questions constitute a fully specified mathematical statement.

\begin{enumerate}
\item Does there exist a field \(F\) such that
\[
\operatorname{char}(F)=2
\qquad\text{and}\qquad
\widetilde u(F)=9?
\]
\item For every field \(F\) in the intended domain, does \(u(F)<\infty\) imply that every anisotropic nonsingular quadratic form \(q\colon V\to F\) satisfying \(\dim_F V=u(F)\) is universal? Explicitly, the requested conclusion is
\[
\forall a\in F^\times\;\exists v\in V\quad q(v)=a.
\]
\end{enumerate}

The second question does not explicitly repeat the characteristic hypothesis. Its intended domain could be fields of characteristic \(2\), continuing the first question, or all fields. These alternatives are not identified here. No positivity restriction on \(u(F)\), exclusion of zero-dimensional forms, or convention for an empty defining supremum is supplied; none is silently added.
\subsection{Short English statement}
The first question asks whether a field of characteristic \(2\) can have \(\widetilde u(F)=9\). The second asks whether every anisotropic nonsingular quadratic form whose dimension equals a finite \(u(F)\) represents every nonzero field element. The invariant definitions and the second question's characteristic scope remain unresolved in the supplied record.
\subsection{Sources}
\begin{itemize}
\item Local registry: \texttt{data/unsolvedmath/problems.json}, permanent UnsolvedMath ID \(30002359\), source problem code \texttt{catalogue entry}, title ``The Value Nine and Universality for Quadratic-Form Invariants.'' The supplied record assigns proposed year \(2013\) and provides no individual proposer.
\item Supplied report citation: \emph{Quadratic Forms and Linear Algebraic Groups} (2014), Oberwolfach Reports, pp.\ 1819--1859, DOI \texttt{10.4171/owr/2013/31}. \url{https://doi.org/10.4171/owr/2013/31}. This is the supplied authoritative source attribution; its contents have not been independently consulted for this draft.
\item Supplied literature source: ``Oberwolfach report on quadratic-form questions.'' \url{https://oa.tib.eu/renate/server/api/core/bitstreams/75e13110-0aa1-46a6-8c43-3bf00c863bc2/content}.
\item Supplied literature source: ``Universal quadratic forms over semi-global fields.'' \url{https://arxiv.org/abs/2309.00689}.
\item Literature assessment dated exactly \texttt{2026-08-21}: the supplied curation labels the questions open, states that existence with \(\widetilde u(F)=9\) remains unresolved, and describes the universality question for arbitrary finite \(u\)-invariant as unsettled, with later work treating important special classes only. This is supplied dated curation, not an independently verified literature assessment. The background explicitly requests independent verification of time-sensitive status and cited claims.
\item The supplied statement status is \texttt{corrected\_verified}; its accompanying verification note says that the statement was checked against the official report and that the clean statement restores source hypotheses and notation. That verification claim is preserved as supplied provenance, not as verification performed here.
\end{itemize}
\end{document}
